\documentclass[12pt]{article}

\usepackage[a4paper, margin=1in]{geometry}
\usepackage{setspace}
\usepackage{enumitem}
\usepackage{titlesec}

\setstretch{1.2}

\title{\textbf{Lecture 23: Predicting the Impossible: Tail Bounds and System Reliability}}
\author{Course: CSE400}
\date{April 02, 2026}

\begin{document}

\maketitle

\section*{Case Study: Amazon Black Friday Sale 2026}

\section*{1. Lecture Information}

\begin{itemize}[leftmargin=0.6in]
    \item \textbf{Course:} CSE400
    \item \textbf{Lecture Number:} 23
    \item \textbf{Title:} Predicting the Impossible: Tail Bounds and System Reliability
    \item \textbf{Case Study:} Amazon Black Friday Sale 2026
    \item \textbf{Topic Focus:} Tails and Bounds
    \item \textbf{Instructor:} Dhaval Patel, PhD
    \item \textbf{Date:} April 02, 2026
\end{itemize}

\section*{2. Amazon Black Friday Sale 2026}

The Amazon Black Friday Sale 2026 is expected to officially begin on Friday, November 27, 2026.

The lecture states that official confirmation from Amazon usually arrives closer to the actual sale date.

The lecture also mentions that historical patterns and expert forecasts from the following sources suggest the expected timeline:

\begin{itemize}
    \item 91mobiles
    \item Bajaj Finserv
\end{itemize}

This sale is used as a case study to discuss uncertainty in computing systems.

\section*{3. Connecting the Sale with Uncertainty and Computing}

The lecture connects the Amazon Black Friday Sale 2026 with uncertainty and computing.

During a large-scale sale event:

\begin{itemize}
    \item User traffic can increase significantly.
    \item The number of requests reaching the system can become unpredictable.
    \item Servers may face sudden spikes in load.
    \item Reliability of the system becomes important.
    \item There is uncertainty in predicting how large the demand surge can become.
\end{itemize}

The lecture uses tail bounds as tools for reasoning about such uncertainty.

\section*{4. Tail Bounds Introduced in the Lecture}

The lecture introduces three important tail bounds:

\begin{itemize}
    \item Markov’s Inequality
    \item Chebyshev’s Inequality
    \item Chernoff Bound
\end{itemize}

These methods are presented as techniques for studying uncertainty and reliability in systems.

\section*{5. Markov’s Inequality}

Markov’s Inequality is introduced as one of the tail bounds.

It is used when we want to bound the probability that a random variable becomes very large.

The lecture includes Markov’s Inequality as a basic probabilistic tool for analyzing uncertain events in computing systems.

Possible application in the Amazon sale context:

\begin{itemize}
    \item Estimating the probability that the number of requests becomes extremely large.
    \item Estimating the probability that system load exceeds a threshold.
\end{itemize}

\section*{6. Chebyshev’s Inequality}

Chebyshev’s Inequality is introduced after Markov’s Inequality.

This inequality is used to reason about how far a random variable can move away from its average value.

The lecture includes Chebyshev’s Inequality as a stronger tool when information about variability is available.

Possible application in the Amazon sale context:

\begin{itemize}
    \item Estimating how far the number of customer requests may deviate from the expected number.
    \item Studying how server demand may vary around the average load.
\end{itemize}

\section*{7. Chernoff Bound}

Chernoff Bound is introduced as another important tail bound.

The lecture presents it as a tool for obtaining stronger probability bounds for large deviations.

Possible application in the Amazon sale context:

\begin{itemize}
    \item Bounding the probability that request traffic becomes much larger than expected.
    \item Estimating the probability of severe overload during the sale.
    \item Studying reliability of large-scale systems under extreme demand.
\end{itemize}

\section*{8. System Reliability}

The lecture links tail bounds with system reliability.

System reliability refers to the ability of a computing system to continue functioning correctly under uncertain conditions.

In the context of the Amazon Black Friday Sale 2026:

\begin{itemize}
    \item Reliability is important because traffic may be very high.
    \item Sudden demand spikes can affect system performance.
    \item Probabilistic bounds help estimate rare but important events.
    \item Tail bounds help in predicting the chance of overload.
    \item Such analysis can help in planning server capacity and handling uncertainty.
\end{itemize}

\section*{9. Summary of the Lecture}

The lecture focuses on the use of tail bounds for reasoning about uncertainty in large-scale computing systems.

The Amazon Black Friday Sale 2026 is used as a motivating example because it involves unpredictable traffic and possible overload conditions.

The three main probabilistic tools introduced are:

\begin{itemize}
    \item Markov’s Inequality
    \item Chebyshev’s Inequality
    \item Chernoff Bound
\end{itemize}

These tools are presented as methods for studying:

\begin{itemize}
    \item Rare events
    \item Large deviations
    \item System overload
    \item Reliability of computing systems
    \item Uncertainty in demand during large online sales
\end{itemize}

\end{document}